\chapter{Measurement Procedure}

This chapter describes a typical measurement workflow.

\section{Preparation}

\begin{enumerate}
  \item Verify all mechanical connections and ensure that all axes move freely.
  \item Power on the system and connect the STM32 and SDRs to the host computer.
  \item Launch the GUI and connect to the STM32 serial port.
\end{enumerate}

\section{Homing and Zeroing}

\begin{enumerate}
  \item Home each axis using the GUI (or manually move to a known reference position).
  \item Confirm that position readouts match expected zero positions.
\end{enumerate}

\section{Configuring a Sweep}

\begin{enumerate}
  \item Set the Azimuth start/stop angles and step size.
  \item Set the Elevation range (if used) and step size.
  \item Configure SDR parameters (center frequency, gain, sample rate).
\end{enumerate}

\section{Running a Measurement}

\begin{enumerate}
  \item Start the sweep from the GUI.
  \item Monitor progress and verify that the AUT moves as expected.
  \item Upon completion, save the measured data (e.g., CSV file).
\end{enumerate}

\section{Data Visualization}

\begin{itemize}
  \item Use built-in plotting in the GUI for quick inspection.
  \item Optionally, export data for use in external tools (MATLAB, Python scripts, etc.).
\end{itemize}
